\section{Introduction}\label{sec:introduction}
Scientific workflows describe the data flow between computational tasks~\cite{Winery}.
With the shift to cloud infrastructures, the question where each task runs determines both the runtime and the cost of a workflow~\cite{Binz2009}.
\lipsum[1]

The contribution of this paper is threefold: (i)~a scheduling heuristic that takes data locality into account, (ii)~an implementation on top of a standard workflow engine, and (iii)~an evaluation with three benchmark workflows (Section~\ref{sec:evaluation}).
\lipsum[2-3]

\section{Background}\label{sec:background}
\lipsum[4]
Figure~\ref{fig:architecture} shows the overall architecture: the scheduler sits between the workflow engine and the resource manager.

\begin{figure}
  \centering
  \includegraphics[width=.6\columnwidth]{example-image-a}
  \caption{Architecture of the scheduling component.}
  \label{fig:architecture}
\end{figure}

The expected transfer time $t$ for a data set of size $s$ over a link with bandwidth $b$ and latency $l$ is
\begin{equation}
  t = l + \frac{s}{b},
  \label{eq:transfer}
\end{equation}
which the heuristic uses to rank candidate hosts (see Equation~\ref{eq:transfer}).
\lipsum[5-6]

\section{Approach}\label{sec:approach}
\lipsum[7]
The core of the heuristic fits into a few lines of pseudo code:
\begin{verbatim}
for task in ready_tasks:
    hosts = rank_by_transfer_time(task)
    host  = first_free(hosts)
    schedule(task, host)
\end{verbatim}
\lipsum[8-9]

\section{Evaluation}\label{sec:evaluation}
Table~\ref{tab:results} lists the makespan of the three benchmark workflows for both strategies.
\begin{table}
  \caption{Makespan in seconds (mean of ten runs).}
  \label{tab:results}
  \centering
  \begin{tabular}{lrrr}
    \hline
    Workflow & FCFS & Locality & Improvement \\
    \hline
    Montage & 412 & 318 & 23\,\% \\
    CyberShake & 1\,205 & 1\,014 & 16\,\% \\
    Epigenomics & 2\,870 & 2\,544 & 11\,\% \\
    \hline
  \end{tabular}
\end{table}
\lipsum[10-12]

\section{Related Work}\label{sec:related}
\lipsum[13-14]
<% const repo = { acmart: "acm-enhanced", ieee: "ieee-enhanced", lncs: "lncs-enhanced" }[documentclass]; -%>
Further material is available at \url{https://github.com/latextemplates/<%= repo %>}.

\section{Conclusion}\label{sec:conclusion}
\lipsum[15]

\section*{Acknowledgment}
The authors thank the anonymous reviewers for their helpful comments.
